\documentclass{article}
\usepackage[T1]{fontenc}
\usepackage{lmodern}
\usepackage{graphicx}
\usepackage{amsmath,amssymb}
\usepackage[a4paper,margin=2cm]{geometry}
\usepackage[absolute,overlay]{textpos}
\setlength{\TPHorizModule}{1mm}
\setlength{\TPVertModule}{1mm}
\usepackage[indonesian]{babel}
\usepackage{hyperref}
\setlength{\emergencystretch}{2em}
% unit: Halaman 35

\begin{document}

% === Halaman 35 (llm) ===

\begin{itemize}
    \item $CFB$ $s$-bit mengenkripsi plainteks berukuran $s$ bit setiap kalinya, $s \le b$ ($b =$ ukuran blok dalam satuan bit).
    \item Dengan kata lain, $CFB$ $s$-bit memperlakukan $cipher$ blok sama seperti pada $cipher$ alir.
    \item Mode $CFB$ membutuhkan sebuah antrian (\textit{shift register}) yang berukuran sama dengan ukuran blok masukan ($b$ bit).
    \item Tinjau mode $CFB$ 8-bit yang bekerja pada blok berukuran 64-bit pada gambar berikut ($s = 8, b = 64$):
\end{itemize}

\end{document}
